\documentclass[10.5pt,a4paper]{article}
\usepackage{fontspec}
\usepackage{geometry}
\usepackage{xcolor}
\usepackage{tikz}
\usetikzlibrary{positioning,shadows}
\usepackage{tcolorbox}
\usepackage{listings}
\usepackage{tabularx}
\usepackage{booktabs}
\usepackage{fancyhdr}
\usepackage{titlesec}
\usepackage{hyperref}
\usepackage{enumitem}
\usepackage{microtype}
\usepackage{multicol}
\geometry{margin=19mm,top=22mm,bottom=20mm}
\definecolor{ink}{HTML}{172235}\definecolor{navy}{HTML}{12304A}\definecolor{teal}{HTML}{00A6A6}\definecolor{mint}{HTML}{E8F7F5}\definecolor{gold}{HTML}{F2B84B}\definecolor{coral}{HTML}{E76F51}\definecolor{paper}{HTML}{F6F9FC}\definecolor{codebg}{HTML}{111C2E}
\setmainfont{Arial}\setsansfont{Arial}\setmonofont{Courier New}
\hypersetup{colorlinks=true,linkcolor=teal,urlcolor=teal,pdftitle={LEXICO Language Documentation},pdfauthor={LEXICO Project}}
\pagestyle{fancy}\fancyhf{}\fancyhead[L]{\sffamily\small\color{navy}LEXICO}\fancyhead[R]{\sffamily\small\color{teal}Language Documentation}\fancyfoot[C]{\sffamily\small\thepage}
\titleformat{\section}{\Large\sffamily\bfseries\color{navy}}{\thesection}{.6em}{}[\color{teal}\titlerule]
\titleformat{\subsection}{\large\sffamily\bfseries\color{navy}}{\thesubsection}{.5em}{}
\titleformat{\subsubsection}{\normalsize\sffamily\bfseries\color{teal}}{\thesubsubsection}{.5em}{}
\lstdefinelanguage{Lexico}{morekeywords={create,set,value,values,print,convert,append,sort,ascendantly,descendantly,count,check,position,extract,if,then,otherwise,when,while,do,for,from,to,step,repeat,until,attempt,up,times,on,failure,define,routine,returns,giveback,emit,outcome,of,with,args,in,take,user,input,message,class,include,inherits,profile,using,override,open,make,file,write,read,line,lines,clear,close,title,into,load,array,table,mesh,matrix,rows,cols,row,column,size,length,sized,at,and,or,not,plus,minus,times,divided,by,mod,true,false,int,float,char,string,bool},sensitive=true,morecomment=[l]{\#},morestring=[b]",morestring=[b]'}
\lstset{language=Lexico,basicstyle=\scriptsize\ttfamily\color{white},keywordstyle=\color{gold}\bfseries,commentstyle=\color{gray},stringstyle=\color{mint},backgroundcolor=\color{codebg},frame=single,rulecolor=\color{teal},breaklines=true,columns=fullflexible,keepspaces=true,showstringspaces=false,aboveskip=5pt,belowskip=5pt}
\newtcolorbox{tipbox}[1]{colback=mint,colframe=teal!70!navy,title={\sffamily\bfseries #1},fonttitle=\small,boxrule=.7pt,arc=2mm,left=3mm,right=3mm,top=2mm,bottom=2mm}
\newtcolorbox{warnbox}[1]{colback=orange!7,colframe=gold,title={\sffamily\bfseries #1},boxrule=.7pt,arc=2mm,left=3mm,right=3mm,top=2mm,bottom=2mm}
\newcommand{\LX}{\textsf{\textbf{LEXICO}}}
\begin{document}
\begin{titlepage}\pagecolor{paper}
\begin{tikzpicture}[remember picture,overlay]
\fill[navy] (current page.north west) rectangle ([yshift=-8.6cm]current page.north east);
\fill[teal] ([yshift=-8.6cm]current page.north west) rectangle ([yshift=-8.84cm]current page.north east);
\node[anchor=north west,text=white,font=\sffamily\bfseries\fontsize{40}{44}\selectfont] at ([xshift=20mm,yshift=-25mm]current page.north west) {LEXICO};
\node[anchor=north west,text=mint,font=\sffamily\fontsize{18}{23}\selectfont] at ([xshift=21mm,yshift=-53mm]current page.north west) {Official Language Documentation};
\node[anchor=north west,text=white,text width=15cm,font=\sffamily\large] at ([xshift=21mm,yshift=-70mm]current page.north west) {A practical guide to writing readable, typed programs with word-based syntax};
\end{tikzpicture}
\vspace*{10cm}
\begin{tipbox}{Learn the language, not the machinery}
LEXICO is a controlled, compiled programming language. This guide teaches the source language: declarations, expressions, conditions, loops, routines, collections, files, modules, and classes. The compiler's internal architecture is intentionally kept out of the way.
\end{tipbox}
\vfill\sffamily\color{navy}\large Version 1.0 \textbullet\ 2026\\[2mm]\color{ink}\normalsize Readable words. Explicit types. Deterministic programs.
\end{titlepage}\nopagecolor
\tableofcontents\newpage
\section{Welcome to LEXICO}
LEXICO is a statically typed, word-based language. Its expressions use words such as \texttt{plus}, \texttt{divided by}, and \texttt{is greater than}, while its blocks use visible \texttt{>} markers. The syntax is intentionally controlled: programs should read naturally, but every accepted form is precise.
\begin{lstlisting}
create int price set value 20;
create int quantity set value 3;
create int total set value price times quantity;
print "Total:", total;
\end{lstlisting}
The example declares three integers, calculates a product, and prints two expressions. A comma in a print statement separates output expressions; LEXICO prints a space between them and a newline at the end.
\begin{tipbox}{The central idea}
Write code that describes intent directly, while retaining explicit types, explicit control flow, and ordinary compiler checks.
\end{tipbox}
\section{Getting started}
\subsection{Requirements and build}
The repository requires a C++17 compiler, CMake 3.20+, Flex, Bison, and LLVM with the C API. `curl` is needed only for optional AI-assisted recovery. From the repository root:
\begin{lstlisting}[language={}]
cmake -S . -B build-cmake
cmake --build build-cmake --target lexico
\end{lstlisting}
The portable helper performs the same CMake workflow:
\begin{lstlisting}[language={}]
python build.py
\end{lstlisting}
\subsection{Run a program}
LEXICO source files use the `.lx` extension:
\begin{lstlisting}[language={}]
lexico hello.lx
lexico --autofix hello.lx
lexico --dump-expanded hello.lx
\end{lstlisting}
`--autofix` enables deterministic local repairs and optional AI parse-error recovery. `--dump-expanded` displays the source after module loading has been expanded.
\subsection{Your first program}
Create `hello.lx`:
\begin{lstlisting}
create string name set value "LEXICO";
print "Hello,", name;
\end{lstlisting}
`create string name` declares a string. `set value` initializes it. The `print` statement evaluates both expressions from left to right.
\section{Language basics}
\subsection{Statements and comments}
Ordinary statements end in `;`. A `\#` begins a comment:
\begin{lstlisting}
create int answer set value 42; # a comment
print answer;
\end{lstlisting}
Keywords are lowercase and case-sensitive. Use `create`, not `Create`, and `string`, not `String`.
\subsection{Blocks and indentation markers}
A block header is followed by lines beginning with one or more `>` characters:
\begin{lstlisting}
if answer is greater than 0 then
> print "positive";
> if answer is equal to 42 then
>> print "the answer";
print "outside";
\end{lstlisting}
One marker enters a block. Two markers enter a nested block. A line without a marker returns to the outer level. Statements inside the block still require semicolons.
\subsection{Names}
Use descriptive identifiers such as `item_count`, `average_price`, and `is_ready`. Avoid reserved keywords as names. The language uses lowercase keywords; identifiers are normally written in lowercase or snake case for readability.
\section{Variables and assignment}
\subsection{Declarations}
\begin{lstlisting}
create int age set value 21;
create float temperature set value 21.5;
create string city set value "Monastir";
create bool online set value true;
\end{lstlisting}
Several variables of the same type can be initialized together:
\begin{lstlisting}
create int x, y, z set value 10, 20, 30;
\end{lstlisting}
The number of names and values must match.
\subsection{Defaults}
A declaration without an initializer uses a typed default:
\begin{lstlisting}
create int count;
create float ratio;
create char initial;
create string message;
create bool ready;
\end{lstlisting}
The defaults are `0`, `0.0`, the null character, the empty string, and `false`.
\subsection{Assignment}
\begin{lstlisting}
create int score set value 10;
set score to score plus 5;
print score;
\end{lstlisting}
The target must already be declared, and the expression must be compatible with its type.
\subsection{Compound assignment}
A bare identifier followed by an arithmetic operator updates that identifier:
\begin{lstlisting}
create int counter set value 0;
counter plus 1;
counter times 2;
\end{lstlisting}
These mean `set counter to counter plus 1;` and `set counter to counter times 2;`. The right side remains a complete expression.
\section{Data types and literals}
\begin{tabularx}{\textwidth}{l l X}\toprule Type & Examples & Use\\\midrule
`int` & `0`, `42`, `-7` & Whole-number calculations\\
`float` & `3.14` & Fractional values\\
`char` & `'A'` & One character\\
`string` & `"hello"` & Text\\
`bool` & `true`, `false` & Logical state\\\bottomrule\end{tabularx}
\subsection{Integers, floats, characters, and booleans}
\begin{lstlisting}
create int width set value 12;
create float price set value 12.50;
create char grade set value 'A';
create bool enabled set value true;
print width, price, grade, enabled;
\end{lstlisting}
Booleans print as `true` and `false` and are not arithmetic operands. Float values are double-precision values in the runtime.
\section{Operators and expressions}
\subsection{Arithmetic}
\begin{tabularx}{\textwidth}{l X}\toprule Word & Meaning\\\midrule
`plus` & addition; string concatenation\\
`minus` & subtraction; first substring removal\\
`times` & multiplication\\
`divided by` & division\\
`mod` & integer remainder\\\bottomrule\end{tabularx}
\begin{lstlisting}
create int a set value 17;
create int b set value 5;
print a plus b;
print a minus b;
print a times b;
print a divided by b;
print a mod b;
\end{lstlisting}
Precedence is parentheses, then `times`/`divided by`/`mod`, then `plus`/`minus`:
\begin{lstlisting}
create int result set value 2 plus 3 times 4;
create int grouped set value (2 plus 3) times 4;
print result, grouped;
\end{lstlisting}
The results are `14` and `20`.
\subsection{Comparisons and logic}
Conditions use `is equal to`, `is not equal to`, `is greater than`, `is less than`, `is greater or equal to`, and `is less or equal to`. Combine conditions with `not`, `and`, and `or`.
\begin{lstlisting}
if (age is greater or equal to 18) and not (blocked is equal to true) then
> print "allowed";
\end{lstlisting}
\section{Strings and text}
\subsection{Concatenation and removal}
\begin{lstlisting}
create string first set value "Lex";
create string second set value "ico";
print first plus second;
print "banana" minus "na";
\end{lstlisting}
`plus` concatenates. `minus` removes the first occurrence of the right-hand substring; if it is absent, the original string remains.
\subsection{Index and length}
\begin{lstlisting}
create string word set value "LEXICO";
print word at 1;
print length of word;
print length of "hello";
\end{lstlisting}
Indexes are zero-based. An invalid index is a runtime error.
\subsection{Extraction}
\begin{lstlisting}
create string text set value "one-two-three";
print extract "two" from text;
print extract from 0 to 3 in text;
\end{lstlisting}
Extraction expressions operate on string sources and use expression values for their search or range arguments.
\section{Input and output}
\subsection{Printing}
\begin{lstlisting}
create int count set value 3;
print "count:", count, count plus 1;
\end{lstlisting}
Arguments are evaluated left-to-right, separated by spaces, followed by a newline.
\subsection{User input}
Input is part of a declaration or assignment:
\begin{lstlisting}
create int age set value take user input with message "Age: ";
set age to take user input with message "Enter age again: ";
\end{lstlisting}
Message parts are comma-separated expressions:
\begin{lstlisting}
create int limit set value 10;
create int value set value take user input with message "Enter a value below ", limit, ": ";
\end{lstlisting}
Boolean input is integer-like: zero becomes `false`; non-zero becomes `true`.
\section{Conditions}
\subsection{If and otherwise}
\begin{lstlisting}
create int score set value 82;
if score is greater or equal to 90 then
> print "A";
otherwise when score is greater or equal to 75 then
> print "B";
otherwise when score is greater or equal to 50 then
> print "C";
otherwise
> print "F";
\end{lstlisting}
Branches are tested in order; the first matching branch runs. Conditions may be nested with deeper `>` markers.
\section{Loops}
\subsection{While}
\begin{lstlisting}
create int i set value 0;
while i is less than 3 do
> print i;
> set i to i plus 1;
\end{lstlisting}
The condition is checked before every iteration.
\subsection{For}
\begin{lstlisting}
for i from 1 to 5 do
> print i;

for i from 0 to 10 step plus 2 do
> print i;

for i from 5 to 1 step minus 1 do
> print i;
\end{lstlisting}
Bounds are inclusive. The iterator is implicit and loop-local. The step must not be zero.
\subsection{Repeat and attempt}
\begin{lstlisting}
create int attempts set value 0;
repeat
> set attempts to attempts plus 1;
until attempts is greater or equal to 3 then
\end{lstlisting}
`repeat` runs its body before checking the condition. A bounded retry loop is written:
\begin{lstlisting}
attempt up to 3 while value is less than 10 do
> set value to value plus 3;
on failure
> print "limit reached";
\end{lstlisting}
The failure block is optional and runs when the limit is reached while the condition remains true.
\section{Routines}
\subsection{Define and return}
\begin{lstlisting}
define routine add with args in take int left, int right returns int do
> giveback left plus right;
\end{lstlisting}
Global routines require an explicit return type. Use `giveback` to return an expression.
\subsection{Call routines}
\begin{lstlisting}
create int total set value outcome of add with args in take 4, 6;
print total;
emit add with args in take 1, 2;
\end{lstlisting}
`outcome of` is an expression. `emit` calls a routine and discards the result. No-argument calls omit the argument clause:
\begin{lstlisting}
define routine answer returns int do
> giveback 42;
print outcome of answer;
\end{lstlisting}
The compiler checks routine existence, argument count, types, and return compatibility.
\section{Meshes and collections}
\subsection{Typed and dynamic meshes}
\begin{lstlisting}
create int mesh nums sized 5;
create int mesh values set values 12, 32, 12, 54, 25;
create mesh mixed set values 12, "hello", 'x', 3.14, true;
\end{lstlisting}
Typed meshes enforce an element type. Dynamic meshes accept mixed scalar values with runtime type tags.
\subsection{Access, append, and sort}
\begin{lstlisting}
print values at 1;
set values value at 1 to 99;
append 40 to values;
append 5 to values at 0;
sort values ascendantly;
\end{lstlisting}
Indexes are zero-based. Typed collections support sorting for supported scalar types.
\subsection{Count, check, and position}
\begin{lstlisting}
print count 12 in values;
print check 99 in values;
print position of 99 in values;
\end{lstlisting}
These return a count, a boolean, and the first matching index or `-1`.
\subsection{Collection loops}
\begin{lstlisting}
for i in values from 0 to size of values minus 1 do
> print values at i;
\end{lstlisting}
\section{Matrices}
\begin{lstlisting}
create int matrix mesh grid set values 1, 2, 3 / 4, 5, 6;
print grid at 0,1;
set grid at 0,1 to 99;
print size of grid;
print row length of grid;
print column length of grid;
\end{lstlisting}
Matrix initializer rows must have equal lengths. Matrix coordinates are zero-based. Matrix loops use two implicit iterators:
\begin{lstlisting}
for i j in grid i from 0 to 1 and j from 0 to 2 do
> set i plus j at i,j;
\end{lstlisting}
\section{File operations}
File statements are valid only inside an active file block:
\begin{lstlisting}
create string line set value "";
create table rows with size 0;
open "notes.txt" file make do
> clear file;
> write "alpha" in file;
> write "beta" in file;
> read line 1 in file into line;
> print "Read:", line;
> read file into rows;
> print rows;
\end{lstlisting}
Use `file do` for an existing file and `file make do` to create a missing file. Single-line reads target a string; multi-line reads target a table. Other forms include `clear line`, range clearing, line replacement, `set file title to`, and `close file`.
\section{Modules}
Load another source file with:
\begin{lstlisting}
load "lib/crypto.lx";
\end{lstlisting}
Relative paths are resolved from the file containing the load statement. Duplicate loads are ignored and circular loads are rejected.
\section{Classes and profiles}
\begin{lstlisting}
define class Shape include
> create string label;
> profile base do
>> set label to "shape";
> routine show do
>> print label;

define class Circle inherits Shape include
> override create string label;
> create int radius;
> profile default do
>> set label to "circle";
>> set radius to 2;
> override routine show do
>> print label, radius;

create Circle c using default;
emit c_show;
\end{lstlisting}
Classes support typed properties, profiles, inheritance, property overrides, routines, and routine overrides. A profile contains property assignments or input actions. Class routines may omit `returns`; routines with values may declare a return type.
\section{Conversions and built-ins}
\subsection{Conversions}
\begin{lstlisting}
create int code set value 65;
convert code to char;
convert code to string;
\end{lstlisting}
The expression form is also available:
\begin{lstlisting}
create string text set value "42";
create int number set value convert text to int;
\end{lstlisting}
Invalid conversions produce runtime errors. Supported built-in expressions include:
\begin{tabularx}{\textwidth}{l X}\toprule Expression & Result\\\midrule
`size of name` & Collection size\\
`row length of name` & Matrix row count\\
`column length of name` & Matrix column count\\
`length of expression` & String length\\
`count value in name` & Match count\\
`check value in name` & Boolean membership\\
`position of value in name` & First matching index or `-1`\\
`name at index` & Collection/string element\\
`name at row,col` & Matrix cell\\\bottomrule\end{tabularx}
\section{Complete programs}
\subsection{Factorial}
\begin{lstlisting}
create int n set value 5;
create int result set value 1;
create int i set value 1;
while i is less or equal to n do
> set result to result times i;
> set i to i plus 1;
print "factorial:", result;
\end{lstlisting}
This program demonstrates declarations, a loop, arithmetic precedence, and output.
\subsection{Mesh search}
\begin{lstlisting}
create int mesh scores set values 12, 45, 78, 45, 91;
create int target set value 45;
print "count:", count target in scores;
print "present:", check target in scores;
print "first index:", position of target in scores;
for i in scores from 0 to size of scores minus 1 do
> if scores at i is greater than 50 then
>> print "high score", i, scores at i;
\end{lstlisting}
\section{Errors and limitations}
LEXICO reports lexical, parse, semantic, and runtime errors. It has no general `try/catch` construct. `attempt ... on failure` is a bounded retry loop, not exception handling. File operations require an active file block; collection indexes are zero-based; typed collections enforce compatibility; and `step 0` is invalid.
\begin{lstlisting}
# Incorrect
Create int value set value 1
if value > 0 then
> print value;

# Correct
create int value set value 1;
if value is greater than 0 then
> print value;
\end{lstlisting}
\section{Language reference and cheat sheet}
\subsection{Common keywords}
`create`, `set`, `value`, `values`, `print`, `convert`, `append`, `sort`, `count`, `check`, `position`, `extract`, `if`, `then`, `otherwise`, `when`, `while`, `do`, `for`, `from`, `to`, `step`, `repeat`, `until`, `attempt`, `up to`, `times while`, `on failure`, `define`, `routine`, `returns`, `giveback`, `emit`, `outcome`, `class`, `include`, `inherits`, `profile`, `using`, `override`, `open`, `make`, `file`, `write`, `read`, `line`, `lines`, `clear`, `close`, `title`, `into`, `load`, `array`, `table`, `mesh`, `matrix`, `size`, `length`, `at`, `int`, `float`, `char`, `string`, `bool`, `true`, `false`.
\subsection{Quick syntax}
\begin{lstlisting}
create int x set value 10;
set x to x plus 1;
print x;
if x is greater than 10 then
> print "large";
while x is less than 20 do
> set x to x plus 1;
define routine double with args in take int n returns int do
> giveback n times 2;
create int mesh nums set values 3, 1, 2;
print check 2 in nums;
open "notes.txt" file make do
> write "hello" in file;
load "library.lx";
\end{lstlisting}
\begin{tipbox}{Programming style}
Keep block markers aligned, use descriptive names, add parentheses when grouping matters, and prefer one clear operation per statement. LEXICO is readable when the source reflects the program's intent.
\end{tipbox}
\end{document}
